\chapter{The C2 boundary: one uniform five-root proof}
\label{component:p9f1235511111}
\noindent\textbf{Component }\path{p9f1235511111}\par
{\small\noindent Source: \path{CORE/prior/inherited/c2_unified.tex}\par}
{\small\noindent SHA256: \path{9f1235511111e57d1ccceec8476e65f1973266a17f78189ee10c9e3bf015e2f8}\par}
{\small\noindent Inventory use: Only exact source/complete parent, centered packet differences, whole scalar/both normal parity and monic division, current elimination, reference Gram computation and complete state/sign join; its unit-N theorem is not imported on the radial domain.\par}
\medskip
\section{The C2 boundary: one uniform five-root proof}
\label{p9f1235511111:sec:c2-unified}

The argument below proves the entire original parent for one explicit
C2 source domain and every independent zero-zero test.  Three source
parameters may approach their boundary at arbitrary relative rates.
Exact contact, parity, and a retained current replace polynomial tables.

\subsection{The unchanged source and its domain}
Let $K_c,K_t$ be independent real closed extensions of $\mathbb R$.
In the coefficient field take $c>0$ infinitesimal and $a,r,g,j>0$ bounded.
Define
\[
 z=c^2,\quad k=1+z,\quad t=1+zr-j,\quad
 B=1-r+g+a^2zt,\quad b=B/k,\quad A_c=1+zr,
\]
and define
\begin{align*}
\Delta={}&(1+z)zb^2+2z(1+z)br+(1+z)(1+a^2z)b\\
 &+z^2(1-a^2)r^2+2z(1+a^2z)r+1+a^2z,\\
N={}&(1+z)(b+r)(b+a^2(1-t))-a^2(1+a^2z)(1-r)t.
\end{align*}
Take $d=1$, $\lambda=\beta=0$, and
\begin{align}
e_1&=\frac{(b+r)A_c(A_c+a^2zt)-r\Delta}{2r\Delta},&
e_2&=\frac{(1+a^2z)(1-r)A_c(A_c+a^2zt)}{2g\Delta},\notag\\
e_3&=\frac{A_cN}{2a^2t\Delta},&
e_4&=\frac{a^2A_c(b+r)(A_c+a^2zt)}{2(b-a^2t)\Delta},\notag\\
e_5&=\frac{(1+z)(b+r)(b-a^2zr)(A_c+a^2zt)}{2a^2j\Delta},&
\alpha&=\frac{z^2(b+r)N}{2(1-z(b-a^2t))\Delta}.
\label{p9f1235511111:eq:c2u-source}
\end{align}
Require the following unit conditions:
\begin{equation}
\begin{gathered}
c>0\text{ is infinitesimal},\\
r,g,j>0\text{ are bounded},\\
a,\quad 1-r,\quad t,\quad b-a^2t,\quad N
\text{ are positive bounded units}.
\end{gathered}\label{p9f1235511111:eq:c2u-hyp}
\end{equation}
Any of $r,g,j$ may have zero or positive residue; no relative-rate
comparison is assumed.

The quantities $b,\Delta$ are positive bounded units.  The source formulas
show that
\begin{equation}
U_1=re_1,\qquad U_2=ge_2,\qquad U_5=je_5,
\qquad e_3,e_4,\qquad \alpha/c^4
\label{p9f1235511111:eq:c2u-units}
\end{equation}
are positive bounded units.  For example, after setting $c=0$ their
first three values are
\[
\frac{b_0(1-r)}{2(b_0+1)},\qquad
\frac{1-r}{2(b_0+1)},\qquad
\frac{b_0(b_0+r)}{2a^2(b_0+1)},\qquad b_0=1-r+g.
\]
Consequently $e_1,e_2,e_5$ have positive bounded reciprocals, whether or
not the coefficients themselves are bounded.  All six event coefficients
are positive.

\begin{theorem}[Uniform C2 unit-gap zero-zero theorem]
\label{p9f1235511111:thm:c2u}
Assume \eqref{p9f1235511111:eq:c2u-hyp}.  For every independent testing pair
$u,v\in K_t$ that is infinitesimal or zero, and every independent real
or complex projective state, the entire parent
\eqref{p9f1235511111:eq:c2u-parent} is a finite sum of squares in
$K_c\otimes_{\mathbb R}K_t$.
\end{theorem}

\subsection{Full parent, contact, and retained squares}
Write $D_u=1+u^2,D_v=1+v^2$, and
\[
\begin{gathered}
R_2=vW-X+uvY-uZ,\qquad R_3=uW-Y,\\
R_5=uW+uvX-Y-vZ,\qquad R_\alpha=vY-Z,\qquad R_\beta=X+uZ.
\end{gathered}
\]
The full homogeneous parent is
\begin{align}
F={}&4e_1D_uD_vX^2+4e_2R_2^2+4e_3D_vR_3^2
 +4e_4D_uD_vY^2+4e_5R_5^2+4\alpha D_uR_\alpha^2\notag\\
 &+2(R_5^2+D_uR_\alpha^2+D_vR_\beta^2-D_uD_vZ^2).
\label{p9f1235511111:eq:c2u-parent}
\end{align}
Its exact coefficient-field stationary zero is
\begin{equation}
(u_c,v_c;W_c,X_c,Y_c,Z_c)
=\left(c,ac;\frac ba,-cr,c\left(\frac ba-at\right),1\right).
\label{p9f1235511111:eq:c2u-contact}
\end{equation}
Substitution of \eqref{p9f1235511111:eq:c2u-source} in the four state derivatives and
two angle derivatives of \eqref{p9f1235511111:eq:c2u-parent} gives zero. These are six
rational identities: one can check them before substituting for $b$ by
using $g=(1+c^2)b-(1-r+a^2c^2t)$ and $j=1+c^2r-t$.
The four state identities imply $F=0$ by homogeneity. Thus this contact
statement is a finite direct calculation, without an asymptotic premise.

First work in the chart $Z=1$, $W\ge0$ bounded. Center three event amplitudes:
\[
 H_1=X+ru,\qquad H_2=R_2-gu,\qquad H_5=R_5+jv.
\]
Put $n_x=H_1,n_y=H_5$. Then, exactly,
\begin{align}
X&=-ru+n_x,&
Y&=uW-(1-j+ru^2)v+uvn_x-n_y,\label{p9f1235511111:eq:c2u-normal-sub}\\
H_2&=G-(1-u^2v^2)n_x-uvn_y,&
G&=D_uvW-[1-r+g+(1-j+ru^2)v^2]u.\label{p9f1235511111:eq:c2u-current}
\end{align}
Define the \emph{entire} remainder
\begin{equation}
\mathcal R=F-4e_1D_uD_vH_1^2-4e_2H_2^2-4e_5H_5^2.
\label{p9f1235511111:eq:c2u-remainder}
\end{equation}
It has bounded coefficient-field coefficients. Indeed the three packet
differences are
\begin{gather*}
-8re_1D_uD_vuH_1+4r^2e_1D_uD_vu^2,\\
8ge_2uH_2+4g^2e_2u^2,\qquad
-8je_5vH_5+4j^2e_5v^2,
\end{gather*}
and \eqref{p9f1235511111:eq:c2u-units} bounds their coefficients. All remaining
packets already have bounded coefficients. Quadratic state degree gives
the exhaustive expansion
\begin{equation}
\mathcal R=f+L_xn_x+L_yn_y
 +Q_{xx}n_x^2+Q_{xy}n_xn_y+Q_{yy}n_y^2,
\label{p9f1235511111:eq:c2u-normal-expansion}
\end{equation}
where all six coefficients are polynomials in $u,v,W$ with bounded
coefficient-field coefficients. There is no higher-normal remainder.

\subsection{Exact parity and Taylor division}
The three soft roots will be
\begin{equation}
 R=au-v,\qquad S=u(W-W_c),\qquad T=u^2-c^2.
\label{p9f1235511111:eq:c2u-soft}
\end{equation}
We prove, with bounded polynomial coefficients,
\begin{equation}
 f\in(R,S,T)^2,\qquad L_x,L_y\in(R,S,T).
\label{p9f1235511111:eq:c2u-membership}
\end{equation}
The parent and its centered substitutions are invariant under
$(u,v,n_x,n_y)\mapsto(-u,-v,-n_x,-n_y)$ with $W$ fixed.
Thus $f(u,au,W)=f_0(u^2,W)$ is even in $u$, whereas
$f_v(u,au,W)$ and $L_i(u,au,W)$ are odd. The polynomial $f_0(\xi,W)$
has degree at most two in $W$. At $u=v=n_x=n_y=0$, the parent is
$4\alpha$, independent of $W$. Hence the nonconstant $W$ coefficients
of $f_0$ are divisible by $\xi$.

At \eqref{p9f1235511111:eq:c2u-contact} every $H_i$ vanishes. Consequently $f$ and
all its first derivatives vanish there, as do both $L_i$. Write
\[
 f_0(\xi,W)=A(\xi)(W-W_c)^2+B_1(\xi)(W-W_c)+C_1(\xi).
\]
Then $\xi\mid A,B_1$, $B_1(c^2)=0$, and
$C_1(c^2)=C_1'(c^2)=0$. Since $c\ne0$, there are polynomials with
bounded coefficients such that
\[
 A=\xi A_1,\qquad B_1=\xi(\xi-c^2)B_2,\qquad
 C_1=(\xi-c^2)^2C_2.
\]
All divisions here are by a monomial or a monic polynomial.
Synthetic division uses only powers of the bounded coefficient $c^2$;
in particular no quotient contains $1/c$. It follows that
\[
 f(u,au,W)=A_1S^2+uB_2ST+C_2T^2.
\]
Next write $f_v(u,au,W)=u\ell(u^2,W)$. Contact stationarity gives
$\ell(c^2,W_c)=0$. Subtracting the value at $W_c$ and dividing that
value by $\xi-c^2$ gives
\[
 f_v(u,au,W)=S\ell_1+uT\ell_2
\]
with bounded coefficients. Exact finite Taylor expansion in $v=au-R$
proves the scalar assertion of \eqref{p9f1235511111:eq:c2u-membership}. The same
odd-parity argument applies to each $L_i(u,au,W)$, and finite division
by $v-au$ proves the normal assertions. This accounts for all scalar,
first-normal, and quadratic-normal terms in the original parent.

\subsection{Eliminate one soft root by a pure tensor unit}
Put $K=1-j+ru^2$ and
\[
Q(\xi)=aW_c-a^2(1-j)-a^2r(\xi+c^2).
\]
The source definition of $b$ gives the exact identity
\begin{equation}
G=-[D_uW-Ku(v+au)]R+aD_uS+uTQ(u^2).
\label{p9f1235511111:eq:c2u-current-soft}
\end{equation}
Indeed $(\xi-c^2)Q(\xi)$ equals
$a(1+\xi)W_c-(1-r+g)-a^2(1-j+r\xi)\xi$.
Together with \eqref{p9f1235511111:eq:c2u-current}, this yields
\begin{equation}
S=\frac{H_2+(1-u^2v^2)H_1+uvH_5
 +[D_uW-Ku(v+au)]R-uTQ(u^2)}{aD_u}.
\label{p9f1235511111:eq:c2u-eliminate-S}
\end{equation}
The denominator is a product of a pure coefficient unit and a pure
testing unit.  Thus this is an allowed exact substitution in the tensor
ring.  Take the five roots
\begin{equation}
K_1=\sqrt{e_1}H_1,\quad K_2=\sqrt{e_2}H_2,\quad
K_5=\sqrt{e_5}H_5,\quad R,\quad T.
\label{p9f1235511111:eq:c2u-five-roots}
\end{equation}
All inverse square roots used to substitute for the $H_i$ are bounded
pure coefficient quantities.  Adding the three retained packets back
therefore gives a bounded exact Gram matrix for the entire parent on
\eqref{p9f1235511111:eq:c2u-five-roots}.

\subsection{A uniform positive residue calculation}
Here is a calculation using only small quadratic forms.  It proves
strictness even when any of the first three source coefficients diverges.
Set $c=0$ formally in the bounded source formulas, and put
\[
\begin{gathered}
b_0=1-r+g,\qquad t_0=1-j,\qquad w_0=b_0/a,\\
N_0=(b_0+r)(b_0+a^2(1-t_0))-a^2(1-r)t_0.
\end{gathered}
\]
Define pure coefficient quantities
\begin{align*}
n&=\frac{b_0(1-r)}{2r(b_0+1)},&
m&=\frac{1-r}{2g(b_0+1)},&
C&=\frac{N_0}{2a^2t_0(b_0+1)},\\
A&=\frac{a^2(b_0+r)}{2(b_0-a^2t_0)(b_0+1)},&
D&=\frac{b_0(b_0+r)}{2a^2j(b_0+1)}.
\end{align*}
The letters $n,m,C,A,D$ in this calculation denote these reference
weights, not additional assumptions or altered source coefficients.
The quantities $C,A,w_0$ are positive bounded units; $n,m,D$ and their
reciprocals are positive, with the reciprocals bounded.

The quadratic angular and hard-state part of the parent at $c=0$ is
\begin{align}
F_2={}&4nX^2+4m(vW-X-u)^2+4C(uW-Y)^2
 +4AY^2+4D(uW-Y-v)^2\notag\\
&+2\bigl((uW-Y)^2-2uvW+X^2+2uX+u^2\bigr).
\label{p9f1235511111:eq:c2u-F2}
\end{align}
Write
\[
\begin{gathered}
h_x=4n+4m+2,\qquad h_y=4C+4A+4D+2,\\
p=\frac{(4C+4D+2)4A}{h_y},\quad
q=\frac{4m(4n+2)}{h_x}.
\end{gathered}
\]
Completing the two hard squares and using the source identities gives
\begin{align}
F_2={}&h_x\left[n_x+\frac{4m}{h_x}(WR-aS)\right]^2
 +h_y\left[-n_y+\frac{4A}{h_y}(S+w_0R/a)\right]^2\notag\\
&+p(S+w_0R/a)^2+q(WR-aS)^2.
\label{p9f1235511111:eq:c2u-F2-factor}
\end{align}
Here the linearized centered substitutions are
$X=-ru+n_x$, $Y=uW-t_0v-n_y$,
$R=au-v$, and $S=u(W-w_0)$.  This identity can be checked by two
square completions and the following five elementary rational identities:
\begin{gather*}
\frac{(4m+2)4n}{h_x}=qa^2w_0^2,\qquad
\frac{4D(4C+4A+2)}{h_y}=p w_0^2/a^2,\\
4m+4D+2-\frac{(4m+2)4m}{h_x}
 -\frac{(4C+4D+2)4D}{h_y}
 =w_0(p/a+qa),\\
(4m+2)-4maw_0=rh_x,\qquad
\frac{4A}{h_y}w_0+\left(\frac{4D}{h_y}-t_0\right)a=0.
\end{gather*}
Their denominators are explicit pure coefficient factors, and direct
substitution of the five displayed weights proves each identity.

At residue level the exact current relation says
$H_2=-WR+aS-n_x$.  Substituting
$S=(H_2+n_x+WR)/a$ in \eqref{p9f1235511111:eq:c2u-F2-factor} cancels the
hard-completion denominators and gives the particularly simple form
\begin{equation}
(4n+2)n_x^2+4mH_2^2+(4C+4D+2)n_y^2
 +4A\left(\frac{(W+w_0)R+n_x+H_2}{a}-n_y\right)^2.
\label{p9f1235511111:eq:c2u-pre-scaled}
\end{equation}
Scale the reference roots by $\sqrt n,\sqrt m,\sqrt D$ and denote
the resulting abstract Gram coordinates by $K_1,K_2,K_5$.  The reference
scaled Gram is
\begin{align}
\mathcal E={}&(4+2/n)K_1^2+4K_2^2
 +\bigl(4+(4C+2)/D\bigr)K_5^2\notag\\
&+\frac{4A}{a^2}\left[
(W+w_0)R+\frac{K_1}{\sqrt n}+\frac{K_2}{\sqrt m}
-\frac{aK_5}{\sqrt D}\right]^2.
\label{p9f1235511111:eq:c2u-reference}
\end{align}
Every coefficient of this Gram matrix is bounded, including when
$r,g,j$ have zero residue.  Its residue is strictly positive definite:
the first three square coefficients are at least $4$, and after those
roots vanish the last square has strictly positive coefficient on $R$
because $A,a,w_0$ are positive units and $W\geq0$.

We justify carefully that this is the residue of the exact Gram just
constructed.  First $re_1-rn$, $ge_2-gm$, and $je_5-jD$ are
infinitesimal, as are $e_3-C$, $e_4-A$, and $W_c-w_0$.
Moreover $e_1/n=U_1/(rn)$, $e_2/m=U_2/(gm)$, and
$e_5/D=U_5/(jD)$ are positive bounded units with residue $1$.
Thus the reference root scaling and the actual scaling in
\eqref{p9f1235511111:eq:c2u-five-roots} induce the same residue coordinates.
After retention and scaling, every coefficient involving the possibly
unbounded packets is a polynomial combination of bounded quantities
such as
\[
U_1,\quad rU_1,\quad \frac1{\sqrt{e_1}}=
\frac{\sqrt r}{\sqrt{U_1}},\quad
\frac{U_1}{\sqrt{e_1}}=\sqrt r\sqrt{U_1},
\]
and the corresponding quantities for $g,j$.  All square-root
denominators here have positive residue.  Their residues thus agree
with those obtained using $rn,gm,jD$.  Monic division, followed by
\eqref{p9f1235511111:eq:c2u-eliminate-S}, also commutes with residue, since its
only further denominator is $aD_u$.  Hence the four-root block is
exactly the residue of \eqref{p9f1235511111:eq:c2u-reference}.

The radial pivot has an equally short exact computation.  On the scalar
graph $v=au,W=W_c$ the remainder is
$(u^2-c^2)^2C_2(u^2)$.  At $u=0$ it equals $4\alpha$; hence
$C_2(0)=4\alpha/c^4$.  The latter is a positive bounded unit by
\eqref{p9f1235511111:eq:c2u-units}.  All radial cross coefficients contain
an explicit factor $u$, by the parity division above; the elimination
\eqref{p9f1235511111:eq:c2u-eliminate-S} adds only further such factors.
Thus the entire five-root residue is
\begin{equation}
\overline{\mathcal E}
 +4\overline{\alpha/c^4}\,T^2\succ0.
\label{p9f1235511111:eq:c2u-strict-residue}
\end{equation}
The strict-residue Gram lemma proves the full bounded $Z=1,W\geq0$
chart, without any restriction on the size of $X,Y$ or of the roots.

\subsection{The other state chart and endpoint dispatch}
For $W=1$, $Z$ infinitesimal, retain the $\alpha$ packet and the three
original squares in $X,R_2,R_5$. The exact identities
\[
Y=uvX-R_5+u-vZ,\qquad R_3=-uvX+R_5+vZ,
\]
\begin{equation}
v=\frac{R_2+(1-u^2v^2)X+uD_vZ+uvR_5}{D_u}
\label{p9f1235511111:eq:c2u-wunit-root}
\end{equation}
express the entire bounded remainder as a bounded quadratic form on
$X,R_2,R_5,u$. For clarity, this does not require a division by any
root: $D_u$ is pure testing. The exact signed-core completion
\[
\begin{split}
H={}&2\left[C_*^2+D_vQ_*^2+D_uv^2Y^2
 +\frac{D_uu^2v^2(1-u^2v^2)}{D_v}Z^2\right]-4D_uuvZ,\\
C_*={}&u+uvR_\beta-Y,\qquad
Q_*=R_\beta-\frac{D_uuv^2}{D_v}Z
\end{split}
\]
shows quadraticity directly after \eqref{p9f1235511111:eq:c2u-wunit-root}: every
displayed amplitude is linear in the four roots, and the last indefinite
term is $u$ times the same linear expression for $v$. All coefficients
are bounded. Adding the three retained packets and using roots
$K_1=\sqrt{e_1}X$, $K_2=\sqrt{e_2}R_2$,
$K_5=\sqrt{e_5}R_5$, and $u$, the residue is the residue of
\[
(4+2/e_1)K_1^2+4K_2^2+
\bigl(4+(4e_3+2)/e_5\bigr)K_5^2
+4e_4\left(u-\frac{K_5}{\sqrt{e_5}}\right)^2.
\]
It is strictly positive because $e_4$ is a positive bounded unit.  This
proves the whole chart, including arbitrary $X,Y$: only the Gram
coefficients, not the roots, needed boundedness.

Finally the testing-only involutions
\[
(u,W,X,Y,Z)\mapsto(-u,-W,-X,Y,Z),\qquad
(v,W,X,Y,Z)\mapsto(-v,-W,X,-Y,Z)
\]
preserve the full parent and allow $u,v\ge0$. Put
$L=1+v^2+u^2v^2$. The $X,Y$ block of half the signed core has eigenvalues
$L\pm uv>0$, and its exact Schur complement is
\begin{equation}
 \frac{uvD_uD_v[uv(W^2+Z^2)-2LWZ]}{L^2-u^2v^2}.
\label{p9f1235511111:eq:c2u-core-schur}
\end{equation}
If $WZ\le0$, the pure-testing core is nonnegative and is a square in
$K_t$; every event is already a tensor square. This includes exact zero
state or angle endpoints. If $WZ>0$, either $W/Z$ is bounded, giving
the first chart, or $Z/W$ is infinitesimal, giving the second. Restoring
the state scale is pure testing multiplication. Splitting a complex
state into real and imaginary parts completes the proof of
Theorem~\ref{p9f1235511111:thm:c2u}.

\paragraph{Scope.}
This theorem includes all eight ordinary or infinitesimal combinations
of $r,g,j$ under \eqref{p9f1235511111:eq:c2u-hyp}; S38 is one special case.
A vanishing residue of $a$, $1-r$, $t$, $b-a^2t$, or $N$ requires
an additional argument.  Complementary testing-angle corners and
entry of every C2 source into this domain are separate obligations.
The proof uses no archived scalar or normal coefficient arrays.
